\section{Matched-information Benchmark}

Table~\ref{tab:benchmark} reports the initial matched-information benchmark. The CNN is clearly predictive and substantially outperforms the linear Logit baseline. However, it does not dominate strong nonlinear tabular models. HistGB-120 attains a higher equal-weighted high-minus-low Sharpe than the CNN, while the validation-selected MLP-4-120 is stronger than either CNN or HistGB in this flattened tabular comparison.

\begin{table}[htbp]
\centering
\caption{Matched-information benchmark}
\label{tab:benchmark}
\begin{tabular}{lrrrr}
\toprule
Model & AUC & Rank IC & EW H-L Sharpe & VW H-L Sharpe \\
\midrule
CNN (E1b-ric) & 0.5190 & 0.0308 & 2.3883 & 1.1098 \\
Logit & 0.5110 & 0.0048 & 1.0954 & 0.2466 \\
HistGB-120 (current chrono panel) & 0.5227 & 0.0256 & 2.5670 & 1.3773 \\
MLP-4-120 (five-seed ensemble) & -- & 0.0410 & 2.8071 & 1.3249 \\
\bottomrule
\end{tabular}
\begin{minipage}{0.92\linewidth}
\footnotesize
\textit{Notes:} All models use the same 120-dimensional 20-day OHLCV information set. The main protocol uses chronological validation and selects hyperparameters by validation Rank IC. The out-of-sample period is 2008--2025.
\end{minipage}
\end{table}


This result changes the interpretation of chart-image models. A CNN can extract useful price-history signals, but image rendering is not necessary for strong return predictability. The stronger conclusion is that nonlinear extraction matters. Once a sufficiently flexible tabular learner is stabilized, much of the apparent chart-image advantage disappears. Section~\ref{sec:temporal_benchmark} pushes this point further by replacing flattened tabular inputs with an explicit temporal numerical architecture.

The flattened MLP benchmark uses the validation-selected four-hidden-layer architecture from the depth-robustness exercise reported below. The CNN benchmark is not presented as a full architecture optimization. We retain the literature-based CNN architecture and conduct targeted market-specific tuning over learning rate, dropout, and early-stopping patience using the 2006--2007 validation period. The 2008--2025 OOS sample is never used for hyperparameter selection. Rank IC is adopted as the primary validation criterion because it is more closely aligned with the paper's cross-sectional ranking objective; in the realized OOS sample, the Rank-IC-selected specification also performs better than the MCC-selected alternative.

\subsection{MLP Stabilization}

A single-seed MLP is sensitive to training randomness. We therefore use five-seed ensembles and select depth by chronological validation Rank IC; the 120D depth comparison selects MLP-4.

\begin{table}[htbp]
\centering
\caption{MLP depth robustness and specification selection}
\label{tab:mlp_stability}
\begin{tabular}{lrr}
\toprule
MLP specification (120D, five-seed ensemble) & Validation Rank IC & OOS Rank IC \\
\midrule
MLP-2 ($128\rightarrow64$) & 0.0705 & 0.0354 \\
MLP-3 ($128\rightarrow64\rightarrow32$) & 0.0691 & 0.0399 \\
MLP-4 ($128\rightarrow64\rightarrow32\rightarrow16$) & 0.0759 & 0.0410 \\
MLP-5 ($128\rightarrow64\rightarrow32\rightarrow16\rightarrow8$) & 0.0737 & 0.0411 \\
\bottomrule
\end{tabular}
\begin{minipage}{0.92\linewidth}
\footnotesize
\textit{Notes:} Architecture selection uses five-seed ensemble validation Rank IC from the 2006--2007 chronological validation period. MLP-4 is therefore the validation-selected flattened nonlinear benchmark; its 180D main result is reported in Table~\ref{tab:shared_specific}.
\end{minipage}
\end{table}

